%---------------------------------------------------------------
% Anhang
%---------------------------------------------------------------
\renewcommand\thesection{\Alph{section}}
\renewcommand\thesubsection{\thesection.\arabic{subsection}}
% VDI-Vorgabe: "Anhang" im Inhaltsverzeichnis und als Überschrift vor den Anhangteilen
\header{Anhang}\label{app}
\counterwithin{table}{section}

%---------------------------------------------------------------------
% Anhang A: 
%---------------------------------------------------------------------
\section{Appendix A}\label{app:A}
%
Tabelle~\ref{tab:tep-stoerungen} listet die im Tennessee-Eastman-Prozess vordefinierten Prozessstörungen auf. Die Störungen IDV(1)--IDV(20) entstammen der ursprünglichen Prozessbeschreibung von \cite{downs.1993} und umfassen Sprünge, zufällige Schwankungen, eine langsame Drift sowie klemmende Ventile; IDV(16)--IDV(20) sind dabei nicht näher spezifiziert. Die Erweiterung um IDV(21)--IDV(28) nach \cite{bathelt.2015} ergänzt zusätzliche, als zufällige Schwankungen modellierte Temperatur- und Druckstörungen der Zulaufströme sowie der Kühlwasser-Umwälzeinheiten. Die Störung IDV(29) stellt die eigene Ergänzung dar. Sie modelliert nach Abschnitt~\ref{sec:tep_drift} eine schleichende Kalibrierdrift des Strippertemperatursensors, die allein den aufgezeichneten Messwert verfälscht und den Prozess selbst unberührt lässt, und dient damit als Grundlage für einen Datenstrom mit driftbehafteten Zusammenhängen.
%
\begin{table}[htbp]
\centering
\footnotesize
\caption{Prozessstörungen des Tennessee-Eastman-Prozesses nach \cite{downs.1993}, erweitert um die zusätzlichen Störungen nach \cite{bathelt.2015}.}
\label{tab:tep-stoerungen}
\begin{tabular}{cll}
\toprule
Nummer & Prozessvariable & Typ \\
\midrule
IDV(1)  & A/C-Zulaufverhältnis, B-Zusammensetzung konstant (Strom 4) & Sprung \\
IDV(2)  & B-Zusammensetzung, A/C-Verhältnis konstant (Strom 4) & Sprung \\
IDV(3)  & D-Zulauftemperatur (Strom 2) & Sprung \\
IDV(4)  & Eintrittstemperatur des Reaktorkühlwassers & Sprung \\
IDV(5)  & Eintrittstemperatur des Kondensatorkühlwassers & Sprung \\
IDV(6)  & Verlust des A-Zulaufs (Strom 1) & Sprung \\
IDV(7)  & Druckverlust im C-Sammler -- verringerte Verfügbarkeit (Strom 4) & Sprung \\
IDV(8)  & Zulaufzusammensetzung von A, B, C (Strom 4) & Zufällige Schwankung \\
IDV(9)  & D-Zulauftemperatur (Strom 2) & Zufällige Schwankung \\
IDV(10) & C-Zulauftemperatur (Strom 4) & Zufällige Schwankung \\
IDV(11) & Eintrittstemperatur des Reaktorkühlwassers & Zufällige Schwankung \\
IDV(12) & Eintrittstemperatur des Kondensatorkühlwassers & Zufällige Schwankung \\
IDV(13) & Reaktionskinetik & Langsame Drift \\
IDV(14) & Reaktorkühlwasserventil & Festklemmen \\
IDV(15) & Kondensatorkühlwasserventil & Festklemmen \\
IDV(16) & Unbekannt & Unbekannt \\
IDV(17) & Unbekannt & Unbekannt \\
IDV(18) & Unbekannt & Unbekannt \\
IDV(19) & Unbekannt & Unbekannt \\
IDV(20) & Unbekannt & Unbekannt \\
\midrule
IDV(21) & A-Zulauftemperatur (Strom 1) & Zufällige Schwankung \\
IDV(22) & E-Zulauftemperatur (Strom 3) & Zufällige Schwankung \\
IDV(23) & A-Zulaufdruck (Strom 1) & Zufällige Schwankung \\
IDV(24) & D-Zulaufdruck (Strom 2) & Zufällige Schwankung \\
IDV(25) & E-Zulaufdruck (Strom 3) & Zufällige Schwankung \\
IDV(26) & A- und C-Zulaufdruck (Strom 4) & Zufällige Schwankung \\
IDV(27) & Druckschwankung in der Kühlwasser-Umwälzeinheit des Reaktors & Zufällige Schwankung \\
IDV(28) & Druckschwankung in der Kühlwasser-Umwälzeinheit des Kondensators & Zufällige Schwankung \\
\textbf{IDV(29)} & \textbf{Aufgezeichnete Strippertemperatur XMEAS(18)} & \textbf{Concept Drift} \\
\bottomrule
\end{tabular}
\end{table}

%
Tabelle~\ref{tab:modes_setpoints} führt die Setpoints auf, mit denen die sechs Betriebsmodi des \ac{TEP} gefahren werden. Sie gehören zum dezentralen Regelkonzept nach \cite{ricker.1996} und beschreiben den jeweils optimalen stationären Betriebspunkt nach \cite{ricker.1995}, übernommen in der Zusammenstellung von \cite{reinartz.2022}. Füllstände, Reaktordruck, Dampfventil und Rührerdrehzahl gehen als Schranken oder feste Vorgaben in die Optimierung ein und bleiben deshalb über alle Modi gleich, ebenso die Stellung des Recycle-Ventils außerhalb der Modi 3 und 6. Die drei letztgenannten Größen sind dabei keine Führungsgrößen, sondern Stellgrößen, die das Regelkonzept konstant hält. Dabei bezeichnet $y_\mathrm{AC}$ den Molanteil von A und C im Reaktorzulauf und $y_\mathrm{A}$ den darauf bezogenen Anteil von A.
%
\begin{table}[htbp]
\centering
\footnotesize
\caption{Setpoints der sechs Betriebsmodi des Tennessee-Eastman-Prozesses nach \cite{ricker.1995, ricker.1996} in der Zusammenstellung von \cite{reinartz.2022}.}
\label{tab:modes_setpoints}
\renewcommand{\arraystretch}{1.2}
\begin{tabular}{l c c c c c c c}
\toprule
& & \textbf{Modus 1} & \textbf{Modus 2} & \textbf{Modus 3} & \textbf{Modus 4} & \textbf{Modus 5} & \textbf{Modus 6} \\
\midrule
Produktionsrate         & \si{\metre\cubed\per\hour} & 22,89  & 22,73  & 18,04  & 36,04  & 23,55  & 20,20  \\
Stripper-Füllstand      & \si{\percent}              & 50     & 50     & 50     & 50     & 50     & 50     \\
Separator-Füllstand     & \si{\percent}              & 50     & 50     & 50     & 50     & 50     & 50     \\
Reaktorfüllstand        & \si{\percent}              & 65     & 65     & 65     & 65     & 65     & 65     \\
Reaktordruck            & \si{\kilo\pascal}          & 2800   & 2800   & 2800   & 2800   & 2800   & 2800   \\
Molanteil G im Produkt  & mol-\%                     & 53,80  & 11,66  & 90,09  & 53,35  & 11,65  & 90,07  \\
$y_\mathrm{A}$          & \si{\percent}              & 63,137 & 64,196 & 62,110 & 61,940 & 64,030 & 61,468 \\
$y_\mathrm{AC}$         & mol-\%                     & 51     & 54,24  & 47,43  & 58,76  & 54,32  & 48,79  \\
Reaktortemperatur       & \si{\celsius}              & 122,9  & 124,2  & 121,9  & 128,2  & 124,6  & 123,0  \\
Stellung Recycle-Ventil & \si{\percent}              & 1      & 1      & 77,62  & 1      & 1      & 71,166 \\
Stellung Dampfventil    & \si{\percent}              & 1      & 1      & 1      & 1      & 1      & 1      \\
Rührerdrehzahl          & \si{\percent}              & 100    & 100    & 100    & 100    & 100    & 100    \\
\bottomrule
\end{tabular}
\end{table}

\section{Appendix B}\label{app:B}
Und der zweite Anhangteil